% !TeX root = seminar-05.tex
% !TeX spellcheck = ru-RU
% !TeX encoding = UTF-8 Unicode
% !BIB program = biber
% LTeX: language=ru-RU
% LTeX: enablePickyRules=true
%
% Copyright (c) 2026 Михаил Михайлов
%
% This work is licensed under the Creative Commons Attribution-NonCommercial-ShareAlike 4.0
% International License. To view a copy of this license, visit:
% http://creativecommons.org/licenses/by-nc-sa/4.0/
%
% You are free to:
%   Share - copy and redistribute the material in any medium or format
%   Adapt - remix, transform, and build upon the material
%
% Under the following terms:
%   Attribution - You must give appropriate credit, provide a link to the license,
%                 and indicate if changes were made.
%   NonCommercial - You may not use the material for commercial purposes.
%   ShareAlike - If you remix, transform, or build upon the material, you must
%                distribute your contributions under the same license as the original.
%
% See the LICENSE file in the repository root for full license text.

\documentclass[seminar, recordsfile={../records.lua}]{practice}

\renewcommand{\SeminarName}{Энтропия, продолжение}
\renewcommand{\SeminarDate}{13 февраля 2026}

\xsimsetup{
	solution/print = false
}

\usepackage{amssymb}
\usepackage{bm}

\DeclareMathOperator{\KL}{KL}
\DeclareMathOperator{\Norm}{Norm}
\newcommand{\h}{\mathrm{h}}

\begin{document}

\begin{exercise}
	Пусть $\xi_1, \ldots, \xi_n$ --- независимые случайные величины, принимающие значения на множестве $\set{1, \ldots, k}$ с вероятностями $p_1, \ldots, p_k$ соответственно. Обозначим за $\eta_{r} = \# i \text{ таких что } \xi_i = r$. Найдите энтропию
	$$
		\H_m(\eta_1, \ldots, \eta_r)
	$$
\end{exercise}
\begin{exercise}
	Пусть $n, k \in \N$, $n \geq k$; Тогда для всех $m > 1$:
	\[
		\log_m\binom{n}{k} \leq n \cdot \H_m(\frac{k}{n})
	\]
\end{exercise}
\begin{exercise}
	Пусть $n \in \N$. Тогда для всех $k \leq \floor{\frac{n}{2}}$ и всех $m > 1$:
	\[
		\sum_{i = 0}^{k}\binom{n}{i} \leq m^{n\H_m(\frac{k}{n})}
	\]
\end{exercise}
\begin{exercise}
	Пусть $g$ --- диффеоморфизм, $\xi$ --- абсолютно непрерывная одномерная с.в., $\eta = g(\xi)$; Выразите $h(\eta)$ через $h(\xi)$.
\end{exercise}
\begin{exercise}
	Найдите $\h(Y)$ для $Y \sim \Exp(\lambda)$.
\end{exercise}
\begin{exercise}
	Найдите $\h(Y)$ для $Y \sim \mathcal{N}(\mu, \lambda)$.
\end{exercise}
\end{document}
